\documentclass[11pt]{article}

\usepackage[margin=1in]{geometry}
\usepackage{amsmath,amssymb}
\usepackage{array}
\usepackage{booktabs}
\usepackage{enumitem}
\usepackage{hyperref}
\usepackage{microtype}
\usepackage[T1]{fontenc}
\usepackage{lmodern}

\hypersetup{
    colorlinks=true,
    linkcolor=blue,
    urlcolor=blue
}

\title{\textit{cs-mail}: Permissioned Communication with Bonded First Contact}
\author{C-SQD}
\date{August 2026}

\begin{document}
\maketitle

\begin{abstract}
\textit{cs-mail} is a communication protocol in which recipients control who may
communicate with them without friction. An accepted sender communicates without
paying per-message fees. A sender who has not been accepted must temporarily
reserve a communication bond consisting of a processing charge and recipient
collateral. Acceptance refunds every decision-open admitted bond for the directed
relationship, fully cancels its unadmitted reservations, and grants standing
communication permission. Rejection transfers the collateral
to the recipient and the processing charge to the recipient provider. Silence
returns the collateral but pays the processing charge to the provider. The design
places cost on unsuccessful attempts to establish a relationship rather than on
communication within an established relationship. This paper explains the
problem, mechanism, incentives, identity model, privacy architecture, and path
from a centralized deployment to an interoperable network. Native message content
is end-to-end encrypted, while pairwise identifiers, compartmentalized services,
and deliberate record expiration protect the communication graph and settlement
history. Exact protocol requirements are defined separately in the
\textit{cs-mail Protocol Specification}.
\end{abstract}

\paragraph{Reading guide.}
Sections 1--3 introduce the protocol through one first contact before adding
formal notation. Later sections add relationship-wide decisions, repeated
attempts, privacy, incentives, settlement, and deployment. The accompanying
protocol specification is normative; this paper explains the design and its
tradeoffs.

\section{The Communication-Permission Problem}

Electronic mail allows anyone who knows an address to impose delivery, filtering,
storage, and attention costs on that address. Spam filters attempt to predict
whether a recipient wants each message, but they operate after the sender has
already exercised an inexpensive right to transmit. Automation and generative
systems continue to reduce the marginal cost of producing plausible unsolicited
messages, while the recipient's time remains scarce.

Charging for every message would address abundance by making ordinary
communication expensive. That is the wrong boundary. Friends, colleagues, and
organizations with an established reason to communicate should not incur a new
financial toll each time they speak. The scarce event is not communication within
a recognized relationship; it is an unknown sender's attempt to establish one.

\textit{cs-mail} makes that boundary explicit. Permission is a persistent,
directed relationship controlled by the recipient. Known and accepted senders
communicate normally. Unknown, rejected, or revoked senders may still attempt
contact, but they must place value behind the attempt. Rejection declines the
present solicitation rather than permanently banning the sender: a later attempt
may become eligible after the applicable backoff. A separate recipient-controlled
block refuses future ordinary contact until the recipient removes it.

\section{A First Contact, End to End}

Alice has not yet been accepted by Bob. Bob's provider quotes a processing
charge of $0.25$ settlement units and collateral of $4.75$ settlement units.
Alice authorizes a temporary reservation of $5.00$ and sends one encrypted
message. The provider admits it only after the encrypted content and its delivery
intent are durably recorded. The reservation does not purchase Bob's attention or
a response; admission makes the contact attempt valid under the protocol.

Bob's relationship decision produces one of three outcomes:

\begin{center}
\small
\begin{tabular}{@{}p{0.13\textwidth}p{0.28\textwidth}p{0.24\textwidth}p{0.24\textwidth}@{}}
\toprule
\textbf{Outcome} & \textbf{Settlement of 5.00} & \textbf{Future permission} & \textbf{Meaning} \\
\midrule
Accept & All 5.00 returns to Alice. & Alice may communicate with Bob without future ordinary bonds. & Bob grants the directed relationship. \\
Reject & 0.25 goes to Bob's provider; 4.75 goes to Bob. & Alice remains unaccepted and may retry only after backoff. & Bob declines this solicitation. \\
Ignore & 0.25 goes to Bob's provider; 4.75 returns to Alice at expiry. & Alice remains unaccepted. & No relationship decision arrived in time. \\
\bottomrule
\end{tabular}
\end{center}

The relationship is directed: Bob accepting Alice does not mean that Alice has
accepted Bob. If Bob accepts while several Alice-to-Bob bonds remain open, the
decision refunds all of them atomically. If Alice remains unaccepted and tries
again, the later attempt may also require backoff and an additional refundable
persistence reserve. Those refinements are introduced only after the basic
first-contact rule below.

\section{The Core Mechanism}

Let sender $i$ attempt to contact recipient $j$. If $j$ has accepted $i$, no bond
is required. Otherwise, the destination side publishes signed contact terms and
the sender reserves

\begin{equation}
    B = C + S,
\end{equation}

where $C$ is a communication-processing charge and $S$ is collateral selected by
the recipient, subject to any minimum established by the recipient provider.
The quote fixes the amounts, settlement unit, duration, and policy version before
the sender commits funds.

The relationship may then be accepted, rejected, or left undecided until the
bond expires:

\begin{equation}
\begin{array}{lll}
\text{Accept} &:& C+S \rightarrow \text{sender},\\[3pt]
\text{Ignore} &:& S \rightarrow \text{sender},\quad
                  C \rightarrow \text{recipient provider},\\[3pt]
\text{Reject} &:& S \rightarrow \text{recipient},\quad
                  C \rightarrow \text{recipient provider}.
\end{array}
\end{equation}

$C$ and $S$ serve different purposes. The processing charge compensates the
destination service when an unsolicited attempt does not become an accepted
relationship. Collateral creates a larger loss when the recipient affirmatively
rejects the attempted relationship. A wanted first contact is costless apart from
the temporary use of liquidity.

Reservation is not admission. The quote supplies a short admission window, and
the sender or an expiry worker may cancel a still-unadmitted reservation. If the
message is not admitted before that window closes, the entire $C+S+L_k$ returns
to the sender, no attempt level advances, and no recipient solicitation appears.
The recipient's decision window begins at canonical admission time. This prevents
abandoned uploads from locking funds or paying a processing charge for a message
the recipient provider never admitted.

\section{Relationships, Not Individual Messages}

Permission is directed. If $j$ accepts $i$, that does not imply that $i$ has
accepted $j$. It is also attached to visible protocol identities: accepting one
address does not automatically accept every address controlled by the same person
or organization.

Most importantly, acceptance and rejection are decisions about the communicative
relationship, not classifications of isolated messages. Several bonded attempts
may be outstanding for the same directed identity pair. Acceptance therefore
performs one atomic relationship transition. It

\begin{itemize}[leftmargin=2em]
    \item establishes permission for future communication;
    \item refunds $C+S$ for every admitted, decision-open bond for the
          relationship;
    \item cancels and fully releases every reserved but unadmitted attempt for the
          relationship;
    \item releases active persistence reserves associated with attempts from the
          accepted public identity; and
    \item cancels the corresponding scheduled expiries.
\end{itemize}

Acceptance does not reset the repeated-attempt state shared by the sender's
private principal and this recipient identity, nor does it release reserves
associated with another public identity controlled by that principal. The
accepted public identity bypasses that anti-abuse state while permission remains
active. Other identities do not inherit its acceptance or forgiveness.

Once the relationship is accepted, an older message cannot subsequently be
rejected. Messages admitted after the transition require no bond. A concurrent
reservation either linearizes before acceptance and is cancelled with a full
release if still unadmitted, or linearizes afterward and is rejected as
unnecessary without retaining value.

Rejection is likewise relationship-wide. It settles all bonds that are pending at
the rejection's canonical ordering point, leaves the sender outside the accepted
set, and preserves the applicable repeated-attempt constraints. It declines the
active solicitation, not all future contact: after the backoff, another eligible
bonded attempt may open a new solicitation. The interface may present accept or
reject beside a message, but the consequence applies to the directed relationship
represented by that message.

A permanent contact prohibition is a different recipient decision. Blocking
closes the active solicitation, settles its decision-open admitted bonds under
the rejection allocation, fully cancels unadmitted reservations, and prevents new
ordinary contact terms or admission until the recipient unblocks the directed
identity pair. Unblocking restores eligibility for future bonded attempts; it
does not accept the relationship, erase the backoff, or reopen settled bonds.

\subsection{One Solicitation, Potentially Several Messages}

The recipient-facing request to establish a relationship is also scoped to the
directed identity pair rather than to one message. The first admitted attempt for
an unaccepted relationship opens a \emph{relationship solicitation episode} and
generates one logical prompt or notification. Its presentation channel may depend
on recipient settings. Later admitted attempts join that active episode. They may
add messages or update a quiet count in the grouped
presentation, but they do not generate another relationship-establishment prompt.

This coalescing does not make the later attempts financially idempotent. Each is
a distinct authenticated attempt with its own content, bond, deadline, and any
required persistence reserve, and admission advances repeated-attempt state. Only
an exact replay of the same command is idempotent. The user-facing solicitation is
coalesced so that repeated messages cannot repeatedly demand a relationship
decision.

Acceptance, rejection, or blocking closes the active episode and applies to the
relationship and every bond open at the decision's canonical time. Rejection
allows a later eligible attempt after backoff; blocking does not. Silence is
different: it is not an \texttt{IgnoreRelationship} command. The recipient simply
makes no decision, and each bond expires under its own deadline. When no admitted
bond in the episode remains decision-open, the episode lapses. A later eligible
attempt, after the applicable backoff, may open a new episode and one new
solicitation. Merely dismissing or hiding the prompt has no permission or
settlement effect.

\subsection{Acceptance Near Expiry}

An explicit recipient decision should not lose merely because a background expiry
worker happens to run at the same moment. At the same time, a settled payment
cannot remain reversible indefinitely. \textit{cs-mail} resolves this tension with
a canonical event order and a bounded finalization rule.

The authoritative service durably timestamps and sequences an acceptance when it
receives it. An acceptance received on or before a bond's decision deadline takes
priority over mechanical expiry. Expiry settlement is finalized only after a
short processing holdback has allowed timely, already-received decisions to enter
the canonical journal. The recipient receives a signed receipt for the decision.

An acceptance received after a bond's decision deadline may still grant future
communication permission and settle every other bond whose decision window
remains open. It does not reopen a bond that has already expired. Thus explicit
intent wins a genuine processing race, while terminal ledger settlement remains
final.

\section{Repeated Attempts}

A sender could otherwise treat repeated losses of $C$ as the cost of searching
for eventual acceptance, particularly because silence returns $S$. The protocol
therefore makes repeated attempts to the same unaccepted recipient progressively
slower and more capital-intensive.

At attempt level $k$, admission may require a backoff interval $\Delta_k$ and a
refundable persistence reserve $L_k$, with

\begin{equation}
    \Delta_{k+1} \geq \Delta_k,
    \qquad
    L_{k+1} \geq L_k.
\end{equation}

The total reservation becomes

\begin{equation}
    R_k = C + S + L_k.
\end{equation}

$L_k$ is not a fee. It is never paid to the recipient or either provider.
Acceptance releases active reserves associated with attempts by the accepted
public identity. It does not release reserves attached to another identity of the
same private principal. Following rejection or expiry, $L_k$ remains unavailable
until its quoted release time and then returns to the sender. Cancellation before
admission releases it immediately. This raises the liquidity required for
persistent contact without creating an incentive for another party to encourage
failed communication.

The attempt level advances on admission rather than settlement, preventing a
sender from opening many low-level attempts while the first remains unresolved.
It is principal-scoped anti-abuse history and is not reset when one public
identity is accepted. A reserved attempt cancelled before admission releases all
funds and does not advance the level. An explicit policy may later decay old
principal-level history, but acceptance of one identity is not that policy.

\section{Identity, Privacy, and Security}

The privacy objective of \textit{cs-mail} extends beyond hiding the words in a
message. A communication system also reveals who attempted to contact whom,
whether the recipient accepted the relationship, how frequently contact was
attempted, and how money moved afterward. Taken together, those facts form a
relationship graph. The graph can be as sensitive as message content, because it
may reveal personal, professional, medical, political, or commercial associations
without revealing a single sentence of correspondence.

The initial C-SQD deployment therefore distinguishes a private principal from the
public protocol identities through which that principal communicates. Alice may
communicate as \texttt{alice@work.example} and
\texttt{alice@personal.example}. These are two public protocol identities: they
are the personas other participants can address, recognize, accept, or reject.
Alice's provider may know that both identities are controlled by the same private
principal. Bob normally sees only the identity Alice uses with him, not the
private principal or the fact that her two identities are connected.

The distinction serves two different purposes. Recipient permission follows the
public protocol identity, because accepting Alice's work identity should not
silently accept every other identity she controls. Repeated-attempt enforcement
follows the private principal with respect to a particular recipient identity,
because changing addresses or operational keys should not erase the history of
unwanted attempts. A privacy-preserving continuity mechanism can carry that
history across a provider change without turning the principal into a public
identifier.

A private principal is the provider-local subject to which control, recovery, and
anti-abuse continuity are attached. It is not necessarily a verified human or
legal person; depending on provider policy, it may represent a person, account,
organization, or managed group. Its identifier is neither a public address nor a
universal identifier shared among providers or recipients. Other domains receive
only the narrowly scoped assertions they need, so acceptance of one identity does
not expose or authorize the principal's other identities.

\subsection{Content and Envelope Privacy}

Native message content should be end-to-end encrypted from the sender's device to
the recipient's device. Content includes the subject, body, inline media, and
attachments. Providers store and transport ciphertext and do not possess the
ordinary decryption keys. The encrypted content remains outside the relationship
state and financial journal, so a ledger compromise cannot become a message
archive and a content-store compromise cannot rewrite settlement.

Delivery still requires a small envelope. The destination must be routable, the
sender must be authenticated, and the recipient provider must be able to verify
that a required bond exists. The envelope should disclose only those facts. A
valid bond does not justify exposing the subject, content, legal identity,
unrelated addresses, account balance, or history of other relationships.

End-to-end encryption protects content from providers and storage compromise, but
it does not protect a device on which the recipient has opened the message. Device
security, authenticated software updates, key recovery, key rotation, and visible
warnings for unexpected key changes are therefore part of the security model.

\subsection{A Private Relationship Graph}

The protocol should not represent a relationship throughout the system by
repeating two email addresses or by hashing those addresses in a predictable way.
It should assign the directed pair an opaque, pairwise-scoped handle. A sender's
relationship with one recipient receives a handle that cannot be correlated with
the handle used for another recipient. The handle means nothing outside the
service and relationship context for which it was created.

An accepted relationship can persist through such a handle or a recipient-issued
capability. Delivery needs to know that the presented authority remains valid; it
does not need a readable history of how the relationship was formed. Revocation
invalidates the authority without publishing an address-book edge. Rejected,
ignored, and pending attempts receive the same confidentiality as accepted ones.

Pairwise handles reduce linkability but do not make traffic analysis impossible.
Timing, message size, network addresses, and distinctive settlement amounts can
still connect otherwise separated events. Providers should limit retention of
these signals, avoid shared tracing identifiers across privacy domains, and use
batching or aggregation where it does not weaken settlement correctness.

\subsection{Separation and Minimal Disclosure}

No ordinary subsystem needs the complete story of a communication. The identity
service needs to know that a verified principal controls a protocol identity, but
not the message content. The content service needs to store encrypted bytes, but
not legal identity evidence or bond ownership. The relationship service needs an
opaque directed-pair handle and its current state. The ledger needs opaque account,
bond, and settlement references together with amounts. The delivery service needs
a destination route and proof that the current permission or reservation rule has
been satisfied. External financial rails need deposits and withdrawals, not a
per-message history.

This separation is meaningful only if the cross-references are protected as
carefully as the underlying records. A central operator that retains every mapping
and every key can reconstruct the graph even when the individual databases look
opaque. C-SQD should therefore separate access roles, encryption keys, operational
logs, and retention schedules so that routine service operation never assembles a
global readable relationship graph. Exceptional reconstruction should be narrow,
audited, and governed by an explicit policy.

\subsection{The Lifecycle of Private Records}

Privacy changes as an attempt moves through the protocol. While a bond is pending,
the system needs enough short-lived linkage to admit the message, prevent replay,
apply backoff, receive a decision, and settle correctly. After finality and any
defined dispute period, most of that event detail is no longer operationally
necessary. The settlement can be reduced to the minimum accounting record that
must remain, while the readable connection among identity, relationship, message,
and ledger records is removed.

An accepted relationship may need to persist for years, but its history need not.
The durable record can be limited to an opaque capability or pair handle, current
permission state, version, and revocation information. It need not preserve every
message that led to acceptance or make the two endpoints discoverable through a
global index.

Encryption alone is not forgetting while the operator retains a usable key and a
mapping back to people. At the end of the applicable operational, dispute, and
legally required retention periods, C-SQD should delete unnecessary records and
destroy the keys that protect expired identifying linkages. This cryptographic
erasure must cover replicas, backups, search indexes, observability systems, and
derived datasets. Anonymous aggregate measurements may remain only when they can
no longer be used to reconstruct a person or relationship.

Historical linkage may also be protected by threshold encryption. For example, a
sender share and a recipient share could both be required to reconstruct a private
relationship archive. This is useful after the protocol no longer requires online
identification of the parties. It is not a substitute for records that must remain
available during settlement, fraud investigation, a dispute, or a legally required
retention period. Any separately retained compliance record should contain the
minimum necessary fields, use separate keys and access controls, expire under a
documented schedule, and never include message content merely because a financial
record must be preserved.

The system can forget only what it controls. A sender and recipient may retain
their own messages and receipts, and an endpoint compromise may reveal what that
endpoint knows. The protocol promise is that C-SQD and participating providers do
not create an unnecessary permanent archive of everyone else's communications and
relationships.

This model creates a substantial governance obligation. The identity repository,
key infrastructure, recovery process, compliance archive, and access audit are
sensitive systems. Their design must minimize disclosure and must also address
participants for whom conventional bank-based or legal-identity verification is
unavailable or inappropriate.

\section{Economic Effects and Counter-Incentives}

The mechanism creates four sender-side costs without taxing accepted
communication:

\begin{enumerate}[leftmargin=2em]
    \item ignored attempts lose $C$;
    \item rejected attempts lose $C+S$;
    \item unsettled bonds consume liquidity; and
    \item repeated attempts encounter increasing backoff and persistence reserves.
\end{enumerate}

A selective sender whose contacts are commonly accepted may therefore operate at
low expected cost. Indiscriminate outreach faces direct losses and working-capital
requirements that scale with unsuccessful volume.

The design must also constrain recipient-side and provider-side incentives. A
recipient can receive $S$ by rejecting, and a recipient provider receives $C$ when
an attempt expires or is rejected. A mature deployment should therefore combine
auditable settlement with clear pricing, prominent and easy acceptance controls,
fraud detection, limits on manipulative solicitation, and dispute procedures.
Providers should not be able to manufacture expiry through hidden decisions or
ambiguous delivery. Signed admission and decision receipts make these behaviors
observable, while governance must define permissible pricing and interface
practices.

A valid bond purchases a protocol-valid opportunity to communicate. It does not
purchase a recipient's attention, a favorable inbox rank, a response, or a right
to bypass safety controls. Malware, unlawful content, excessive volume, and other
policy violations remain independently rejectable at the transport and service
layers.

\section{Ledger and Settlement}

The financial ledger is part of protocol validity because a bonded attempt is not
valid unless the required value has been reserved. Three layers nevertheless
remain distinct:

\begin{equation}
\boxed{
\text{protocol settlement semantics}
\neq
\text{ledger operator}
\neq
\text{banking rails}
}
\end{equation}

Protocol semantics determine who owns reserved value after each event. The ledger
operator maintains an append-only, auditable record of those obligations. Banks,
card processors, and other regulated services fund and redeem ledger balances.
Many message settlements can consequently occur internally without creating a
bank transaction for each message.

Federated providers may clear those internal obligations in aggregate. For
example, if Provider A owes Provider B $10{,}000$ during a clearing period and
Provider B owes Provider A $8{,}000$, an external rail may settle the $2{,}000$
net obligation while both providers retain the individual authenticated protocol
events needed for participant receipts, audit, and dispute handling. Netting
changes how providers discharge obligations; it does not change the ownership
result of any individual protocol transition.

The journal should use double-entry accounting, conserve value in every
transition, and make all commands idempotent. Relationship changes and their
financial consequences must be atomic. A bond and a persistence reserve may each
reach terminal settlement exactly once.

\section{Launch Extension: Express Lanes for Scoped Communication}

Permanent personal acceptance is not appropriate for every relationship. A
merchant, airline, healthcare organization, public agency, or automated service
may need permission to communicate only about a transaction, appointment, account,
or bounded period.

Recipients may therefore issue scoped capabilities defining an authorized sender,
recipient, purpose or message class, validity interval, volume limit, and
revocation conditions. A matching capability permits bond-free delivery without
granting unrestricted standing permission. Capabilities preserve the same central
principle: communication authority originates with the recipient.

Express lanes implement this capability outside the settlement kernel. A
capability-authorized message neither opens nor joins a relationship solicitation
nor advances repeated-attempt state. A recipient block atomically invalidates
previously issued capabilities for the directed identity pair.

\section{Deployment Path}

An initial deployment can be centralized without collapsing the logical roles in
the design. C-SQD may operate identity verification, the canonical event journal,
permission state, bond ledger, delivery infrastructure, and internal clearing,
while external financial providers handle funding and withdrawal.

The first implementation should focus on a deterministic protocol kernel. Given
authenticated commands, explicit time, and prior state, the kernel computes the
same relationship and ledger transitions without performing networking, database
access, or external payments. Unit tests and property-based tests can then verify
relationship-wide settlement, value conservation, replay safety, and expiry races.
The surrounding services should be separated from the beginning into identity,
relationship, encrypted content, ledger, and telemetry domains. Retrofitting that
separation after a global identifier has entered every log and database would not
recover the privacy that was already lost.

SMTP gateways can support conventional senders during migration. Later federation
can allow a sender provider to reserve funds and a recipient provider to verify
the reservation, deliver the message, and produce authenticated decisions.
Providers may net their obligations periodically while preserving the individual
protocol events. Federation requires shared representations for identities,
quotes, messages, bonds, decisions, receipts, and clearing obligations, but it
does not require identical internal databases or banking partners.

Operator-specific choices for identity enrollment, account funding, SMTP payment
challenges, clearing, and staged launch are described separately in the
\textit{cs-mail Deployment and Migration Profile}. They are not requirements of
the protocol itself.

\section{Conclusion}

\textit{cs-mail} treats communication permission as durable protocol state and
places economic friction at the first-contact boundary. Accepted communication is
free of protocol message charges. Unaccepted communication carries a refundable
bond whose settlement reflects whether the recipient accepts or rejects the
relationship or instead makes no decision before individual bonds expire.

The central mechanism is simple, but its implementation depends on precise
semantics. Acceptance and rejection apply to relationships rather than isolated
messages; rejection declines the current solicitation while a distinct block
prohibits future ordinary contact. Relationship decisions atomically settle all
eligible admitted bonds and fully cancel unadmitted reservations. A canonical
ordering rule gives timely acceptance priority over mechanical expiry without
weakening finality. Repeated-attempt controls attach to private principals, while
visible permission and acceptance-side reserve release remain attached to public
identities; accepting one identity does not reset the principal's shared history.
One active solicitation groups repeated attempts so that they cannot repeatedly
demand the same relationship decision.

The same architecture treats privacy as a property of the communication graph as
well as its content. Native content is end-to-end encrypted. Relationships and
attempts use pairwise opaque references, and each service receives only the facts
needed for its role. After settlement and the applicable retention period, the
system removes unnecessary linkage and destroys expired keys rather than keeping
a permanent readable history of who contacted whom.

These rules provide a path from a centralized product to an interoperable
communication network. The accompanying protocol specification turns them into
versioned objects, commands, state transitions, error cases, and testable
invariants.

\end{document}
